\documentclass{newlayout}
%Bitte hier den enstprechenden Ort einsetzen z.B. Braunschweig und die Akademienummer
\Akademie{Ort}{2014}{1}

\usepackage[ngerman, english]{babel}
\usepackage{misc}
\usepackage{multicol}

%\usepackage{amsmath}%wird automatisch durch newlayout.cls geladen
\usepackage{amsfonts}

\usepackage{blindtext}

\usepackage{url}
\def\UrlBreaks{\do\a\do\b\do\c\do\d\do\e\do\f\do\g\do\h\do\i\do\j\do\k\do\l%
\do\m\do\n\do\o\do\p\do\q\do\r\do\s\do\t\do\u\do\v\do\w\do\x\do\y\do\z\do\0%
\do\1\do\2\do\3\do\4\do\5\do\6\do\7\do\8\do\9\do\-\do\_\do\/\do\%}
\urlstyle{same}

% hinzugef�gt, um Fehler 'pdfTeX error (font expansion): auto expansion is only possible with scalable' zu vermeiden
\usepackage{lmodern}
\setkomafont{descriptionlabel}{\normalfont\bfseries}
\addtokomafont{paragraph}{\normalfont}
\usepackage{footnote}
\usepackage[flushmargin,hang,ragged]{footmisc}
\deffootnote{1em}{1em}{%
\textsuperscript{\thefootnotemark\ }
}
%\setlength{\abovedisplayskip}{5pt}
%\setlength{\belowdisplayskip}{5pt}


%%%%%Mathe-Definitionen
\newtheorem{Def}{Definition}
\newtheorem{Sat}{Satz}
\newtheorem{Bew}{Beweis}
\newtheorem{Thm}{Theorem}

\setlength\abovedisplayshortskip{0pt}
\setlength\belowdisplayshortskip{0pt}
\setlength\abovedisplayskip{3pt}
\setlength\belowdisplayskip{3pt}
%%%%Ende Mathe-Definitionen

\begin{document}

 %   \input{titel}
 \setcounter{page}{3}

\setcounter{tocdepth}{1}
 \tableofcontents

   \setcounter{secnumdepth}{1}


\setcounter{page}{7}
\setcounter{chapter}{0}

%Angabe, bis zu welcher Stufe die sections im Text nummeriert werden sollen.
      \settocdepth{2}



\course{1.2}{Lineare Algebra}%%% 
\begin{coursetitle}
  \centerline{Learning from Data aka i++} 
  \bigskip
  \Large \centerline{Mathematische Methoden des digitalen Zeitalters}
  \bigskip
\end{coursetitle}
%\begin{dsafigure}
%\begin{center}
%\includegraphics[width=.9\columnwidth]{Titelbild-fehlt.png}
%\caption{meine Bildunterschrift}
%label{fig:meinbild}
%\end{center}
%\end{dsafigure}



%\documentclass{article}

%\usepackage[ngerman]{babel}
%\usepackage[utf8]{inputenc}
%\usepackage{color}
%\usepackage[a4paper,lmargin={2cm},rmargin={2cm},
%tmargin={2.5cm},bmargin = {2.5cm}]{geometry}
%\usepackage{amssymb}
%\usepackage{amsmath}
%\usepackage{graphicx}
%\usepackage{multicol}
%\usepackage{amsthm}

%\setlength{\columnsep}{1cm}
%\newtheorem{defi}{Definition}
%\newtheorem*{defi*}{Definition}

%\begin{document}
%\begin{multicols*}{2}

\section{Vektoren}
In der linearen Algebra nehmen Vektoren eine zentrale Rolle ein. Ein Vektor wird in der Schule als eine Menge an Pfeilen gelehrt, die parallel, gleichgerichtet und gleich lang sind. Betrachten wir jedoch nicht nur den dreidimensionalen Raum, sondern $\mathbb{R}\textsuperscript{n}$, so ist die Vorstellung eines Pfeils nicht immer m\"oglich . Dies bedingt die Notwendigkeit einer anderen Definition.
\newline
\indent Ein Vektor ist eine Aufz\"ahlung von Objekten und beschreibt eine Verschiebung. Eine solche Aufz\"ahlung entspricht der Definition eines Tupels. Entscheidend bei Tupeln ist die Reihenfolge der Objekte, die auch mehrfach vorkommen k\"onnen.
\begin{equation*}
\text{n-Tupel } (a_1,a_2,\dots,a_n)
\end{equation*}
\subsection{Vektor- \& Unterraum}
Vektoren bilden die Elemente eines Vektorraumes V. Addieren wir zwei Vektoren eines Vektorraumes oder multiplizieren wir sie mit einem Skalar, so ist die Summe bzw. das Produkt ebenfalls ein Element des Vektorraumes.
\begin{align*}
(1)\text{ }u, v \in V,\text{ } u+v \in V  \\
(2)\text{ }u \in V, \lambda \in \mathbb{R} ,\text{ }\lambda u \in V
\end{align*}
Ein Unterraum ist eine Teilmenge eines Vektorraumes. Es gelten f\"ur sie die oben genannten Eigenschaften eines Vektorraumes.
\begin{align*}
U\subseteq V, \text{wenn gilt:}	\\
(1)\text{ }u, v \in U,\text{ } u+v \in U  \\
(2)\text{ }u \in U, \lambda \in \mathbb{R} ,\text{ }\lambda u \in U
\end{align*}
\subsection{Span}
Alle m\"oglichen Vektoren eines Vektorraumes werden durch den sogenannten Span dargestellt. Der Span ist die Menge aller m\"oglichen Linearkombinationen der Basisvektoren.
\begin{align*}
\text{Span }(V_1,\dots,V_k) = \{U \in V : U = \lambda_1v_1+\dots+\lambda_kv_k\} \\
\text{ mit } \lambda_1,\dots,\lambda_k \in \mathbb{R}
\end{align*}
\subsection{Basis \& Dimension}
Die Menge der Basisvektoren wird Basis eines Vektorraumes genannt. Jeder Vektorraum besitzt eine Dimension p, die durch die Anzahl der Basisvektoren bestimmt wird.
\subsubsection*{\"Ubung}
Beweise, dass ein Span immer ein Vektorraum ist.
\begin{proof}
Wir nehmen an: 
\begin{align*}
U \in Span(V_1,\dots,V_k)\\
V \in Span(V_1,\dots,V_k)\\
\\
U = \lambda_1V_1+\dots+\lambda_kV_k\\
V = \alpha_1V_1+\dots+\alpha_kV_k\\
\end{align*}
Additionsregelung bei Vektorr\"aumen (s.(1))
\begin{align*}
	&W_1 = U+V\\
	&=\lambda_1V_1+\dots+\lambda_kV_k+\alpha_1V_1+\dots+\alpha_kV_k\\
	&= (\lambda_1\alpha_1)V_1+\dots+(\lambda_k\alpha_k)V_k
\end{align*}
Multiplikationsregelung bei Vektorr\"aumen (s.(2))
\begin{align*}
	&W_2 = \lambda_1V_1+\dots+\lambda_kV_k\\
	&\lambda W_2= (\lambda \lambda_1)V_1+\dots+(\lambda\lambda_k)V_k
\end{align*}
\end{proof}
\subsection{Lineare Unabh\"angigkeit}
$v_1,\dots,v_k$ mit $v_i \in V$ sind linear unabh\"angig, falls $\lambda_iv_i+\dots+\lambda_kv_k = 0$ nur f\"ur $\lambda_i=\dots=\lambda_k= 0$ gilt. 

%\end{multicols*}
%\end{document}

%TODO \input{Vektoren2}

\section{Lineare Abbildungen}

In diesem Abschnitt werden wir die elementarste Form einer Funktion kennenlernen, die einen Vektorraum auf einen anderen Vektorraum abbildet. Diese Funktionen nennt man \textit{lineare Abbildungen}.
\begin{Def}
	\label{Def:Def_1}
Eine Funktion $f:U \stackrel{}{\rightarrow} V$ ist linear, wenn folgende Bedingungen gelten: \\ 
1.
\begin{equation*}
 f(u+v) = f(u)+f(v) \forall \ u,v \in U
\end{equation*}\\
2.
\begin{equation*}
 f(\lambda v) = \lambda f(v) \ \forall \ v \in U \ und \ \lambda \in \mathbb{R}
\end{equation*}
\end{Def}

Das besondere an solchen Abbildungen ist, dass sie von \textit{Matrizen} der Form $A \in \mathbb{R}^{n \times m}$ definiert werden, wobei $n$ die Dimension von $U$ und $m$ die Dimension von $V$ ist.

\subsection{Matrix}
Eine Matrix ist eine spezielle Anordnung von $n\cdot m$ Zahlen in $n$ Zeilen und $m$ Spalten. Man kann eine Matrix aber auch als eine Zusammenfassung von $m$ Vektoren mit $n$ Komponenten ansehen, dabei bildet ein Vekotor eine Spalte. Die Komponente in der $i$-ten Zeile und $j$-ten Spalte wird als $a_{ij}$ geschrieben und die ganze Spalte mit $a_j$ notiert.\\
\begin{equation*}
A= \left(
   \begin{array}{cccc}
	  a_{11} & a_{12} & ... & a_{1m}\\
		a_{21} & a_{22} & ... & a_{2m}\\
		...   & ...   & ... & ...  \\
		a_{n1} & a_{n2} & ... & a_{nm}
	 \end{array}
	 \right) = (a_1, a_2, ..., a_m)\\
\end{equation*}
\\
Da wir nun wissen was eine Matrix ist, führen wir zwei neue Operationen ein.
\begin{enumerate}
	\item Matrix-Vektor-Multiplikation
	\item Matrix-Matrix-Multiplikation
\end{enumerate}

\subsection{Matrix-Vektor-Multiplikation}
Bei einer Matrix-Vektor-Multiplikation wird eine Matrix $A \in  \mathbb{R}^{n \times m}$ mit einem Vektor $v \in \mathbb{R}^m$ multipliziert. Dabei wird jede Zeile der Matrix mit dem Vektor skalarmultipliziert; deshalb ist es wichtig, dass die Matrix so viele Spalten wie der Vektor Dimensionen hat. Eine Multiplikation zwischen Matrizen und Vektoren, bei denen die Anzahl nicht übereinstimmt, ist nicht definiert.
\begin{align*}
	Av &= 
	\left(
   \begin{array}{cccc}
	  a_{11} & a_{12} & ... & a_{1m}\\
		a_{21} & a_{22} & ... & a_{2m}\\
		...   & ...   & ... & ...  \\
		a_{n1} & a_{n2} & ... & a_{nm}
	 \end{array}
	\right) 
	\times 
	\left(
	 \begin{array}{c}
	  v_1\\
		v_2\\
		...\\
		v_m
	 \end{array}
	\right)\\
	&= \left(
	 \begin{array}{c}
	 a_{11}\times v_1 + a_{12} \times v_2 + ... + a_{1m} \times v_m\\
	 a_{21}\times v_1 + a_{22} \times v_2 + ... + a_{2m} \times v_m\\
		...\\
	 a_{n1}\times v_1 + a_{n2} \times v_2 + ... + a_{nm} \times v_m\\
	 \end{array}
	\right)
\end{align*}
\caption{Matrix mit Vektor multiplizieren}

Allgemein lässt sich für die Matrix-Vektor-Multiplikation sagen
\begin{equation*}
(Av)_i = \sum_{j=1}^{n}a_{ij}v_j 
\end{equation*}

\subsection{Matrix-Matrix-Multiplikation}
Eine Matrix-Matrix-Multiplikation ist nichts anderes als mehrere Matrix-Vektor-Multiplikationen hintereinander. Haben wir zwei Matrizen $A \in \mathbb{R}^{n \times m}$ und $B \in \mathbb{R}^{m \times p}$ und berechnen $AB$ nehmen wir jede Spalte von $B$ als Vektor und multiplizieren ihn mit $A$ wie oben beschrieben. Die Vektoren die man als Ergebnisse erhält, fasst man wieder in einer Matrix zusammen mit der Dimension $\mathbb{R}^{n \times p}$. Allgemein können wir sagen:
\begin{equation*}
	(AB)_{il} = \sum_{n=i}^{m}a_{ij} \times b_{jl}
\end{equation*}

\subsubsection*{Übung}
Beweise, dass eine Komposition $h(x) = f \circ g (x)$ aus den linearen Abbildungen $f$ und $g$, auch eine lineare Abbildung ist.

\begin{proof}
Da $f$ und $g$ linear sind gelten die Bedingungen aus Definition \ref{Def:Def_1}:
\begin{align*}
	h(u+v) &= f \circ g(u+v)\\
	&= f(g(u+v))\\
	&= f(g(u)+g(v))\\
	&= f(g(u)) + f(g(v))\\
	&= f\circ g (u) + f\circ g (v) = h(u) + h(v)
\end{align*}
Des Weiteren gilt:
\begin{align*}
	h(\lambda v) &= f \circ g (\lambda v)\\
	&=f(g(\lambda v))\\
	&= f(\lambda g(v))\\
	&=\lambda f(g(v))\\
	&=\lambda f \circ g (v) = \lambda h(v)
\end{align*}
Da beide Bedingungen aus Definition \ref{Def:Def_1} für $h$ erfüllt sind, ist auch eine Komposition aus zwei linearen Abbildungen linear.
\end{proof}



%\input{Kapitel-2-Kapitel�berschrift}
%\input{Kapitel-3-Kapitel�berschrift}
%usw.

%\nocite{*} % alle referenzen anzeigen, sogar wenn sie nicht im text zitiert sind
\bibliography{lit}{}
\bibliographystyle{plain}
\end{document}
